\documentclass[12pt]{article}
\usepackage[utf8]{inputenc}
\usepackage[T2A]{fontenc}
\usepackage[russian]{babel}
\usepackage{url}

\usepackage{amsmath, amssymb}
\usepackage{physics}
\usepackage{geometry}
\usepackage{graphicx}
\usepackage{setspace}
\geometry{margin=2.5cm}

\title{Physics-Informed Neural Networks:\\
интуитивное и теоретическое введение}
\author{}
\date{}

\begin{document}

\section{Обозначения и соответствие переменных модели и реализации}

\subsection{Независимые переменные}

\begin{itemize}
\item $x,y,z$ --- пространственные координаты точки среды.
\item $t$ --- время.
\item $\Omega \subset \mathbb{R}^3 \times \mathbb{R}$ --- пространственно-временная область.
\end{itemize}

В реализации используется тензор
\[
\texttt{xyzt} \in \mathbb{R}^{N \times 4},
\]
где каждая строка имеет вид $(x,y,z,t)$.

\subsection{Искомые поля}

Нейросеть аппроксимирует поля
\[
\mathcal{N}_\theta(x,y,z,t)
=
\big(u_x,u_y,u_z,T\big).
\]

Где:

\begin{itemize}
\item $\mathbf{u}=(u_x,u_y,u_z)$ --- вектор смещений частиц среды.
\item $T$ --- температурное поле.
\item $\theta$ --- параметры нейросети.
\end{itemize}

В коде:
\begin{verbatim}
pred = model(xyzt)
u = pred[:, 0:3]
T = pred[:, 3:4]
\end{verbatim}

\subsection{Параметры материала}

\begin{itemize}
\item $\rho$ --- плотность среды (кг/м$^3$).
\item $\lambda$ --- первый параметр Ламе.
\item $\mu$ --- модуль сдвига.
\item $k$ --- коэффициент теплопроводности.
\item $C$ --- объёмная теплоёмкость.
\item $\gamma$ --- коэффициент термомеханической связи.
\item $T_0$ --- опорная температура.
\end{itemize}

В коде:
\begin{verbatim}
rho = mat["rho"]
lam = mat["lame_lambda"]
mu  = mat["lame_mu"]
\end{verbatim}

\subsection{Тензор деформаций}

\[
\varepsilon_{ij}
=
\frac{1}{2}
\left(
\frac{\partial u_i}{\partial x_j}
+
\frac{\partial u_j}{\partial x_i}
\right).
\]

\begin{itemize}
\item $\varepsilon_{ij}$ --- тензор малых деформаций.
\item $i,j \in \{x,y,z\}$.
\end{itemize}

В коде:
\begin{verbatim}
eps = strain_tensor_from_xyzt(u, xyzt)
\end{verbatim}

\subsection{След тензора деформаций}

\[
\mathrm{tr}(\varepsilon)
=
\varepsilon_{xx}
+
\varepsilon_{yy}
+
\varepsilon_{zz}.
\]

В коде:
\begin{verbatim}
tr_eps = trace_3x3(eps)
\end{verbatim}

\subsection{Тензор напряжений}

\[
\sigma_{ij}
=
\lambda \delta_{ij} \mathrm{tr}(\varepsilon)
+
2\mu \varepsilon_{ij}
-
\gamma \delta_{ij} T.
\]

Где:

\begin{itemize}
\item $\sigma_{ij}$ --- тензор напряжений.
\item $\delta_{ij}$ --- символ Кронекера.
\end{itemize}

В коде:
\begin{verbatim}
sigma = stress_iso_thermo(eps, T, lam, mu, gamma)
\end{verbatim}

\subsection{Уравнение движения}

\[
\rho \frac{\partial^2 u_i}{\partial t^2}
=
\frac{\partial \sigma_{ij}}{\partial x_j}
+
f_i(x,t).
\]

\begin{itemize}
\item $u_{tt}$ --- ускорение частиц.
\item $\nabla \cdot \sigma$ --- дивергенция тензора напряжений.
\item $f_i(x,t)$ --- внешняя сила (источник).
\end{itemize}

В реализации формируется невязка:
\[
r_{\text{motion}}
=
\rho u_{tt}
-
\nabla \cdot \sigma
-
f.
\]

\subsection{Уравнение теплопереноса}

\[
k \Delta T
-
C \frac{\partial T}{\partial t}
=
\gamma T_0
\frac{\partial}{\partial t}
\mathrm{tr}(\varepsilon).
\]

Где:

\begin{itemize}
\item $\Delta T$ --- оператор Лапласа.
\item $T_t$ --- производная температуры по времени.
\end{itemize}

Невязка:
\[
r_{\text{heat}}
=
k \Delta T
-
C T_t
-
\gamma T_0 \frac{\partial}{\partial t}
\mathrm{tr}(\varepsilon).
\]

\subsection{Источник}

\[
f(x,t)
=
A
\exp
\left(
-\frac{\|x-x_0\|^2}{2\sigma^2}
\right)
w(t)
\hat d.
\]

\begin{itemize}
\item $A$ --- амплитуда.
\item $x_0$ --- координата источника.
\item $\sigma$ --- пространственная ширина источника.
\item $\hat d$ --- единичный вектор направления.
\end{itemize}

Временная часть (волна Рикера):

\[
w(t)
=
(1-2a^2)e^{-a^2},
\quad
a=\pi f_0 (t-t_0).
\]

\subsection{Поток энергии}

\[
\mathbf{v}
=
\frac{\partial \mathbf{u}}{\partial t}
\]

— скорость частиц.

\[
\mathbf{S}
=
-
\boldsymbol{\sigma}
\cdot
\mathbf{v}
\]

— вектор потока механической энергии.

Нормированное направление распространения:

\[
\hat{\mathbf{s}}
=
\frac{\mathbf{S}}{\|\mathbf{S}\|}.
\]

\subsection{Определение направления распространения энергии}

После обучения вычисляется вектор потока энергии
\begin{equation}
\mathbf{S}
=
- \boldsymbol{\sigma} \cdot \mathbf{v},
\qquad
\mathbf{v}
=
\frac{\partial \mathbf{u}}{\partial t}.
\end{equation}

Нормированное направление распространения энергии:
\begin{equation}
\hat{\mathbf{s}}
=
\frac{\mathbf{S}}{\|\mathbf{S}\|}.
\end{equation}

\end{document}